\pg{440} 18. \textbf{O atheismo pratico é possivel.} --- O atheismo \emph{pratico} não é uma depravação da \emph{intelligencia,} mas sim da \emph{vontade.} Ora a nossa vontade está sujeita á depravação, e, para se entregar mais livremente ás desordens, póde deixar de pensar em \textsc{Deus} e de Lhe prestar as homenagens da sua piedade e obediencia. --- Este \emph{atheismo} não só é \emph{possivel,} mas existe realmente\footnote{\emph{Bayle} diz: « Os que não servem a \textsc{Deus}, e que antes o offendem com uma vida cheia de crimes e vicios, são uma especie de atheus, é o atheismo dos libertinos e devassos; porque effectivamente vivem, como se não houvesse \textsc{Deus}, sem medo da sua Justiça, sem gratidão para com a sua Bondade, sem respeito para o seu Nome, sem obediencia ás suas leis, e, ainda que podessem tirar \textsc{Deus} do seu throno para o esmagar sob os seus pés e aniquilar para sempre, não viveriam com mais licença e ousadia. São atheus pelas suas acções; conhecem a \textsc{Deus} com a lingua, mas negam-no com os costumes; fallam como crentes, mas operam como incredulos » (\emph{Pensées div.,} t. IV, pag. 94).}.

19. \textbf{O atheismo theoretico-negativo não é possivel nos homens adultos.} --- A existencia de \textsc{Deus}, embora nos não seja \emph{immediatamente} evidente, todavia descobre-se logo que começamos a applicar os primeiros principios da razão. Por quanto, sendo natural e facillimo deduzir da existencia dos effeitos a existencia da causa respectiva, é claro que, contemplando o universo, não podemos deixar de reconhecer uma causa, que o